% ECONOMICS_PROBLEM: {"id": "D83-459", "title": "Personalization and competition", "jel_code": "D83", "source_id": "OP-0400"}
% FIRST_POSTED: 2026-10-09
% ATTEMPT_STATUS: unresolved
% ENTRY_KIND: research_attempt
% !TEX program = pdflatex
\documentclass[11pt]{article}
\usepackage[T1]{fontenc}
\usepackage[utf8]{inputenc}
\usepackage[margin=27mm]{geometry}
\usepackage{amsmath,amssymb,amsthm,mathtools}
\usepackage[colorlinks=true,linkcolor=blue,citecolor=blue,urlcolor=blue]{hyperref}
\allowdisplaybreaks
\newtheorem{theorem}{Theorem}
\newtheorem{proposition}[theorem]{Proposition}
\newtheorem{lemma}[theorem]{Lemma}
\newtheorem{corollary}[theorem]{Corollary}
\theoremstyle{remark}
\newtheorem{remark}[theorem]{Remark}
\newcommand{\E}{\mathbb E}
\newcommand{\R}{\mathbb R}
\newcommand{\Ind}{\mathbf 1}
\newcommand{\Var}{\operatorname{Var}}
\newcommand{\KL}{\operatorname{KL}}
\newcommand{\CS}{\operatorname{CS}}
\newcommand{\supp}{\operatorname{supp}}
\title{Personalized ranking with optimal limited search: welfare accounting and an equilibrium reversal}
\author{GPT 6 Astra Ultra}
\date{9 October 2026}
\begin{document}
\maketitle
\noindent\textbf{Problem ID:} \texttt{D83-459}. \textbf{Status:} Unresolved; rigorous restricted results.

\section{Target and scope}
The target is the twelve-month consumer-surplus effect of a personalized ranking after sellers adjust prices and quality, compared with the direct effect at fixed suppliers. Farronato \cite{farronato}, Section 1, discusses personalization and pricing responses; Section 1.1 discusses interference and equilibrium adjustment. We prove a search-feasible price decomposition, derive a fully specified two-seller example where a positive fixed-price benefit becomes a negative equilibrium effect, and state what market randomization would identify. These are restricted results, not an estimate for a travel platform.

\section{Search plans, rather than free full information}
Fix a user type $i$ and ranking regime $r$. There is a finite inspection tree whose nodes reveal information about a finite set of sellers. At each node the user may inspect an available item, stop and book an inspected item, or take an outside option. A plan is a decision rule measurable in the information revealed so far. Inspections have specified nonnegative costs. For the next result, the information law, feasible tree, values, qualities, and costs are fixed independently of the price vector $p$. Prices are known at the outset; their only effect is through terminal money payments and which plan is selected. This assumption allows a user to ignore a new price vector when choosing a plan but does not allow use of unobserved values.

Every deterministic feasible plan $\pi$ has an expected gross value net of inspection cost $g_{i\pi}$ and booking probabilities $d_{i\pi j}\geq0$, with $\sum_jd_{i\pi j}\leq1$. Randomization over plans cannot improve on their maximum. Thus
\[
 V_i(r,p)=\max_{\pi\in\Pi_i(r)}\{g_{i\pi}-d_{i\pi}\cdot p\}. \tag{1}
\]
The no-search outside plan belongs to $\Pi_i(r)$ and yields zero. The finiteness assumption can allow many histories and seller states; it is not an assumption that the user learns every value.

\begin{theorem}[Exact price-response decomposition]
Let prices move along $p(t)=p^0+t(p^1-p^0)$, $0\leq t\leq1$, holding regime $r=1$ and the above search primitives fixed. Choose an optimal plan at each $t$ and let $d_i(t)$ be its booking-probability vector. Then
\[
 V_i(1,p^1)-V_i(1,p^0)
   =-\int_0^1d_i(t)\cdot(p^1-p^0)\,dt.                  \tag{2}
\]
Consequently, for any fixed baseline user population with integrable values,
\[
 \Delta\CS=G-\int_0^1\E[d_i(t)]\cdot(p^1-p^0)\,dt,
 \quad G=\E[V_i(1,p^0)-V_i(0,p^0)].                      \tag{3}
\]
Moreover, $|\Delta\CS-G|\leq\|p^1-p^0\|_\infty$. If every price rises, the supplier-price term lies between $-\|p^1-p^0\|_\infty$ and zero.
\end{theorem}
\begin{proof}
Along the price segment (1) is the maximum of finitely many affine functions of $t$. It is continuous and piecewise affine. At any point where a unique affine slope is active, its derivative is $-d_i(t)\cdot(p^1-p^0)$. Multiple plans with different slopes can tie only at finitely many crossing points; plans tying on an interval have the same directional slope. Hence the derivative formula holds almost everywhere, irrespective of tie selection. Integration over finitely many affine pieces proves (2). Adding and subtracting $V_i(1,p^0)$ proves the user-level decomposition. Booking probabilities sum to at most one, so the integrand has absolute value at most $\|p^1-p^0\|_\infty$; this bound justifies Fubini over user types and proves (3) and both bounds.
\end{proof}

If personalized information permits emulating every baseline plan at no greater cost, then $G\geq0$: map a baseline optimizer to its emulating plan and compare values in (1). A change in ranking does not automatically satisfy this condition. It may hide alternatives, change which comparisons are feasible, or worsen some groups' search. The theorem makes no such dominance claim without the emulation assumption. Applied separately to a prespecified baseline group, (3) gives the same exact accounting and bounds.

For quality changes, one robust extension is available. Suppose the two quality vectors share a coupling of the same information tree and histories, and every plan's gross value changes by at most $\varepsilon_i$ under that coupling, while its booking probabilities and search costs stay fixed. For any such plan its objective changes by at most $\varepsilon_i+\|\Delta p\|_\infty$. Taking maxima in both directions proves
\[
 |V_i(r,p^1,q^1)-V_i(r,p^0,q^0)|
       \leq\varepsilon_i+\|\Delta p\|_\infty.            \tag{4}
\]
Without the common-history condition, quality may change information acquisition and (4) needs additional assumptions; bounding realized utility alone does not automatically bound changes in the experiment.

\section{A two-seller price game with an exact sign reversal}
There are two sellers, zero marginal costs, no entry, and fixed qualities. A unit mass of users is split equally into two observed preference types. A user has value $HZ$ for the preferred seller and $LZ$ for the other seller, where $Z$ is uniform on $[0,1]$, independent of type, and $H>L>0$. The value is learned only on inspection. Initial information contains type, prices, and the value distribution, not $Z$.

Each user may inspect only the top-ranked listing, at zero inspection cost, and then either book it or leave. This is a hard one-inspection attention constraint. Lower listings cannot be inspected or booked. Both regimes provide the outside option. The common regime puts either seller first with probability $1/2$ independently of user type, as an exogenous display randomization; the personalized regime always puts the preferred seller first. The common policy is therefore type-independent, even though its displayed top item is randomized. Sellers announce one uniform price each before rankings are realized and maximize expected revenue. The ranking rule does not depend on their prices. Every user optimally inspects, since inspection is free, and books exactly when learned value exceeds price.

\begin{proposition}[Positive direct benefit and negative adjusted benefit]
Suppose $2L<H<3L$. The unique price equilibrium of the common regime has both prices
\[
 p^0=\frac{HL}{H+L},
\]
and that of the personalized regime has both prices $p^1=H/2$. Consumer surplus is
\[
 \CS(0,p^0)=\frac{H^2-HL+L^2}{4(H+L)},\qquad
 \CS(1,p^1)=\frac H8.
\]
The adjusted effect is strictly negative:
\[
 \Delta\CS=-\frac{(H-L)(H-2L)}{8(H+L)}<0,              \tag{5}
\]
whereas $G=\CS(1,p^0)-\CS(0,p^0)>0$.
\end{proposition}
\begin{proof}
A seller under the common regime is inspected by mass $1/4$ of each type. Its demand and revenue are
\[
 D_0(p)=\tfrac14(1-p/H)_++\tfrac14(1-p/L)_+,
 \qquad R_0(p)=pD_0(p),\quad p\geq0.
\]
For $0\leq p\leq L$ this is a strictly concave quadratic with unique maximizer $p^0=HL/(H+L)$ and revenue $HL/[4(H+L)]$. For $L\leq p\leq H$, its maximizer is $H/2$ because $H>2L$, with revenue $H/16$. The first revenue exceeds the second exactly when $H<3L$. Prices above $H$ yield zero revenue. Thus $p^0$ is the unique global optimum. Under personalization, the seller faces only preferred users, of mass $1/2$, and revenue $p(1-p/H)_+/2$ is uniquely maximized at $H/2$. A seller's demand does not depend on its rival's price under this hard search constraint, so these unique choices form the unique equilibria.

For a user inspecting a listing with value uniform on $[0,v]$, expected optimal surplus is $(v-p)_+^2/(2v)$ by direct integration of $(z-p)_+$ over that uniform law. Thus at any common price $p\leq L$,
\[
 \CS(0,p)=\frac{(H-p)^2}{4H}+\frac{(L-p)^2}{4L},\qquad
 \CS(1,p)=\frac{(H-p)^2}{2H}.
\]
Substitution of $p^0$ gives the stated baseline expression, and substitution of $p^1$ gives $H/8$. Subtraction and factorization yield (5). Finally, for each fixed $p<H$, uniform value $HZ$ strictly improves expected positive-part surplus over $LZ$ whenever $p<H$ and $H>L$. Personalization replaces the one-half chance of the inferior listing by the preferred listing, so $G>0$ at $p=p^0$.
\end{proof}

For example, $H=5/2$ and $L=1$ give $p^0=5/7$, $p^1=5/4$, and $\Delta\CS=-3/112$. Direct surplus at baseline prices is $125/196$, versus baseline surplus $19/56$. The price response more than offsets this direct gain. These are exact fractions in a constructed model, not estimated platform effects. No user is granted free knowledge of both sellers, and there is no claim that personalization always raises prices. Allowing users to inspect the rival could restore active substitution and alter every pricing conclusion.

\section{What an experiment must observe}
Market randomization to full twelve-month ranking regimes, with no cross-market spillovers and a specified supplier adjustment rule, identifies the distribution of \emph{observed} market outcomes under each regime. To infer $\CS$, one still needs an identified utility and search model or a valid direct monetary-surplus measurement. Clicks and bookings alone do not identify the unobserved gross values in (1). This limitation remains even under perfect randomization.

Conditional on a calibrated model identifying $V_i$ in each arm, randomization identifies the average of (3) at the selected twelve-month prices and qualities. An additional short-run randomized ranking experiment at baseline suppliers identifies $G$, if the calibrated primitives remain invariant. Their difference then identifies the supplier-response contribution. Individual assignment inside one market generally holds sellers at a mixed-regime response; it identifies $V_i(r,p^{\rm mix})$, not the two market-equilibrium values in (3). The explicit example illustrates this distinction: the seller's demand under a mixed rollout is a different weighted mixture and its maximizing price generally differs from either endpoint.

A useful next proof problem is to replace the hard one-inspection constraint by a sequential costly-search model with substitution, endogenous quality and capacity. Equation (3) still guides that work when its information invariance conditions hold, but a fully estimated booking-platform equilibrium and group-specific empirical welfare remain unresolved.

\section{Completion Estimate}
50\%

This is a post hoc, heuristic estimate for the full catalogue target.
The search feasible surplus decomposition and a fully solved two seller pricing example establish how personalization benefits can reverse after supplier adjustment. Identification of actual search values, competing sequential search, endogenous quality and capacity, market spillovers, and twelve month group specific welfare remain unresolved.

\begin{thebibliography}{9}
\bibitem{farronato} C. Farronato (2025), \emph{Digital Platforms as Information Aggregators: Implications for their Design and Regulation}, conference draft. Sections 1 and 1.1, PDF pp.2--5, inspected 9 October 2026. \url{https://conference.nber.org/confer/2025/DTs25/farronato.pdf}.
\end{thebibliography}

\end{document}
